%! AP Statistics — Unit 3, Lesson 3.7 — Activity KEY
\documentclass[10pt]{article}
\usepackage{apstats-article}
\usepackage{apstats-key}

\begin{document}

\pageheader{Unit 3, Lesson 3.7}{Group Activity KEY: Scope of Inference Stations}
\namepartnerperiod

\vspace{0.1in}

\begin{scenariobox}[Context]{sky}
\small
For each study, your group will complete a full scope-of-inference analysis:
(1) identify the type of randomization used, (2) determine what conclusions are
valid, (3) write a proper conclusion sentence using precise language, and
(4) explain what would need to change to expand the scope of inference.
\end{scenariobox}

\vspace{0.1in}

\textbf{Study A:} A researcher recruits 80 college students who volunteer to
participate. She randomly assigns 40 to a mindfulness app and 40 to a wait-list
control. After 6 weeks, the mindfulness group reports significantly lower stress.

\begin{practicebox}
\small
\begin{enumerate}[leftmargin=1.5em, itemsep=7pt]
  \item Random assignment used? \makebox[3cm][l]{\ans{Yes}} \quad Random selection used? \makebox[3cm][l]{\ans{No (volunteers)}}
  \item Can we conclude causation? \makebox[2cm][l]{\ans{Yes}} \quad Can we generalize to all college
        students? \makebox[2cm][l]{\ans{No}}
  \item Write a proper conclusion sentence:\\
        \ansline{Random assignment lets us conclude the app caused lower stress,}\\[1pt]\ansline{but only for volunteers like these, not all college students.}
  \item What would need to change to allow generalization?\\
        \ansline{Randomly select the participants from all college students}\\[1pt]\ansline{(e.g., from registration lists) instead of using volunteers.}
\end{enumerate}
\end{practicebox}

\vspace{0.1in}

\textbf{Study B:} Researchers randomly select 600 adults from a national census
registry and survey them about their social media use and self-reported happiness.
They find that heavier social media users report lower happiness.

\begin{practicebox}
\small
\begin{enumerate}[leftmargin=1.5em, itemsep=7pt]
  \item Random assignment used? \makebox[3cm][l]{\ans{No}} \quad Random selection used? \makebox[3cm][l]{\ans{Yes}}
  \item Can we conclude social media \emph{causes} lower happiness? \makebox[2cm][l]{\ans{No}}
        Explain:\\
        \ansline{This is an observational study --- no treatment was imposed,}\\[1pt]\ansline{so confounding variables cannot be ruled out.}
  \item Can we generalize to the national adult population? \makebox[3cm][l]{\ans{Yes}}
  \item Name a possible confounding variable: \makebox[9cm][l]{\ans{Depression, loneliness, or age.}}
\end{enumerate}
\end{practicebox}

\vspace{0.1in}

\textbf{Study C:} A biomedical company randomly selects 200 patients from all
patients at a large hospital and randomly assigns them to a new pain drug or a
placebo. Patients taking the new drug report significantly less pain.

\begin{practicebox}
\small
\begin{enumerate}[leftmargin=1.5em, itemsep=7pt]
  \item Random assignment used? \makebox[3cm][l]{\ans{Yes}} \quad Random selection used? \makebox[3cm][l]{\ans{Yes}}
  \item Can we conclude causation? \makebox[3cm][l]{\ans{Yes}} \quad Can we generalize? \makebox[3cm][l]{\ans{Yes}}
  \item To whom can we generalize? Be specific.\\
        \ansline{To all patients at that hospital (the population sampled),}\\[1pt]\ansline{not to all patients everywhere.}
  \item Write a full conclusion sentence addressing both causation and generalization:\\
        \ansline{Because patients were randomly assigned, the new drug caused}\\[1pt]\ansline{less pain; because they were randomly selected, this result}\\[1pt]\ansline{generalizes to all patients at this hospital.}
\end{enumerate}
\end{practicebox}

\vspace{0.12in}

\begin{practicebox}
\small
\textbf{Group Synthesis:} Your group reads this AP free-response excerpt:

\vspace{4pt}
\textit{``An experiment used random assignment but not random selection. The
researchers concluded that the treatment caused the observed improvement and that
the results can be generalized to all adults.''}

\vspace{6pt}
Is this conclusion correct? Identify the error and write a corrected conclusion.\\[6pt]
\ansline{Partly wrong: without random selection we cannot generalize to all adults.}\\[1pt]\ansline{Corrected: random assignment shows the treatment caused the}\\[1pt]\ansline{improvement, but only for subjects like those in the study.}
\end{practicebox}

\end{document}
